\section{Gate cost and hardware viability}\label{sec:hardware}
\begin{table*}[t]\centering\small\setlength{\tabcolsep}{3.5pt}
\caption{\orig{compiled counts} Gate cost of amplitude amplification, standard Grover against $A_q$ (CZ after transpilation to $\{$CZ, $R_z$, $\sqrt X$, $X$, $R_y\}$, no hand optimisation; the objective-threshold oracle is common to both and not counted). Standard iteration: feasibility oracle (automaton compute, phase, uncompute) + diffusion. $A_q$ iteration: $A_q^\dagger$, $A_q$, and the phase flip on all qubits. ``default: Qiskit's ancilla-free multi-controlled gates; ``V-chain: the two all-qubit phase flips use Toffoli V-chains with clean ancillas (extra $N-3$ qubits for $A_q$). Totals = iterations $\times$ cost per iteration (+ one $A_q$).}\label{tab:gen_cost}
\begin{tabular}{lrrrrrrrrrr}\toprule
& & & & & & \multicolumn{2}{c}{total CZ, default} & \multicolumn{2}{c}{total CZ, V-chain}\\\cmidrule(lr){7-8}\cmidrule(lr){9-10}
family & $n$ & $w$ & CZ($A_q$) & CZ(oracle) & $k_{\rm std}\!\to\!k_{A_q}$ & standard & $A_q$ & standard & $A_q$\\\midrule
cardinality & 10 & 5 & 912 & 2\,608 & $25\to17$ & 76\,300 & 129\,908 & 66\,400 & 34\,674\\
knapsack & 10 & 14 & 3\,227 & 9\,732 & $17\to14$ & 172\,992 & 207\,319 & 166\,260 & 96\,271\\
independent set & 10 & 9 & 1\,351 & 4\,912 & $25\to8$ & 133\,900 & 80\,023 & 124\,000 & 24\,407\\
set packing & 10 & 4 & 335 & 1\,660 & $14\to1$ & 29\,456 & 5\,169 & 23\,912 & 1\,143\\
exact cover & 10 & 4 & 276 & 1\,864 & $25\to1$ & 57\,700 & 5\,368 & 47\,800 & 972\\
3-colouring & 10 & 15 & 1\,810 & 7\,598 & $25\to4$ & 201\,050 & 48\,786 & 191\,150 & 17\,058\\
assignment $3\times3$ & 9 & 5 & 358 & 1\,998 & $8\to0$ & 18\,576 & 358 & 16\,320 & 358\\
assignment $4\times4$ & 16 & 9 & 2\,078 & 10\,654 & $201\to4$ & 2\,442\,954 & 108\,310 & 2\,158\,338 & 19\,974\\
\bottomrule\end{tabular}\end{table*}

Iteration counts are not the whole cost: $A_q$ and the reflection about $A_q|0\rangle$ are circuits too. This appendix accounts for the gates, first abstractly and then routed to the \texttt{ibm\_fez} topology, and says which instances could be predicted to run today. Nothing here was executed on a quantum device.

\paragraph{Gate cost per amplification (\cref{tab:gen_cost}).} We also synthesised the feasibility oracle that standard Grover needs (automaton run into fresh registers, phase, uncomputation; verified to flip exactly the feasible strings) and compare totals on equal footing, with the objective-threshold oracle left out because both methods need it. With Qiskit's default ancilla-free multi-controlled gates, the phase flip on all qubits of the $A_q$ iteration dominates (22\,402 of 26\,558 CZ for $4\times4$ assignments) and the totals favour $A_q$ on six of eight rows, by $1.7\times$ (independent set) to $52\times$ ($3\times3$ assignment) and $23\times$ ($4\times4$), but not on the two loose constraints (cardinality $130\,000$ against $76\,000$, knapsack $207\,000$ against $173\,000$). If the two all-qubit phase flips use Toffoli V-chains with clean ancillas (extra qubits), $A_q$ wins on all eight rows, by $1.7\times$ to $108\times$. These are unoptimised counts and the standard side's oracle uses the same unoptimised gates, so the figures show where the iteration savings survive the extra cost of $A_q$ and the reflection: for tight constraints they do, for loose ones they depend on how the multi-controlled phase is synthesised.

\paragraph{Depth across constraint families (compile only).}\label{sec:general-depth} Without executing anything, we compiled $A_q$, one amplification iteration and one Grover-mixer QAOA layer on the \texttt{ibm\_fez} topology (FakeFez snapshot) for seven families (\path{experiments/hardware_depth_multi.py}, \path{experiments/hardware_depth_variants.py}). Settings: optimisation level 3, approximation degree $1$, median of five transpiler seeds, multi-controlled Z gates as Toffoli V-chains; every circuit is legal on the coupling map. With the generic compiler, $A_q$ alone costs $24$--$694$ CZ, and an iteration or QAOA layer $206$--$4388$ CZ (about $0.2$--$3.6$ times the median $T_2$ of the qubits used). Only the exact-cover instance has an estimated success probability near $0.5$ for an iteration (a product of gate errors, not a simulation); for the other families it is $\le0.03$. Routing roughly doubles the CZ count: for the exact-cover iteration, $106$ CZ before routing become $206$--$243$, and for the Dicke iteration $267$ become $505$--$535$.

Three structural changes were tested. (i) For count constraints, a Dicke-state preparation (\path{negsearch/dicke.py}; \cref{rem:anyA}) reduces an iteration for $3$ of $6$ from $1659$ to $529$ CZ ($28\to9$ qubits, estimated success probability $0.001\to0.18$) and $A_q$ alone from $361$ to $118$ CZ; for $4$ of $8$, from $3885$ to $1069$ CZ. (ii) A one-hot automaton register lowers depth on wide automata (assignment $3\times3$: $7508\to4851$) but raises CZ counts on narrow ones (exact cover: $220\to333$), so it is not a general improvement. (iii) For QAOA on count constraints, an XY-ring mixer is much cheaper per layer than the Grover mixer (\cref{tab:gen_xy}). Four attempts gave no usable gain: other multi-controlled synthesis methods ($496$--$574$ CZ against $529$), a layout found by solving the placement as a quadratic assignment problem (\path{negsearch/qap_layout.py}; worse than the default layout on all five circuits tested, for instance $512$ against $227$ CZ for the exact-cover iteration, because the static cost ignores the swaps the router inserts and, without an error term, lands on noisy qubits), many transpiler seeds (about $5$--$10$\,\% fewer CZ), and approximate transpilation. The last is invalid, not merely ineffective: with approximation degree $0.9$ the compiled circuits have $53$ and $61$ CZ before routing but amplify nothing ($P(\mathrm{opt})=0$ in both cases, against $1.0$ and $0.392$), and degree $0.99$ already changes the Dicke result from $0.392$ to $0.659$. We therefore expect the exact-cover and the Dicke-based count instances to be the only ones whose amplification iteration is measurable on current devices.

\begin{table}[t]
\centering\small
\caption{\orig{exact ideal simulation; compiled counts and ESP (model); nothing executed} QAOA on the $3$-of-$6$ cardinality instance ($|F|=20$, one optimum; a uniform feasible sample is optimal with probability $0.05$): XY-ring mixer against Grover mixer, both on the Dicke state. Best $P(\mathrm{opt})$ over $12$ noiseless local optimisations of the angles per depth $p$ (a local optimum, not a global guarantee); CZ and estimated success probability (ESP, product of gate errors) after compilation to the \texttt{ibm\_fez} topology, nothing executed.}\label{tab:gen_xy}
\begin{tabular}{lrrrrrr}
\toprule
 & \multicolumn{3}{c}{XY-ring mixer} & \multicolumn{3}{c}{Grover mixer}\\
\cmidrule(lr){2-4}\cmidrule(lr){5-7}
$p$ & $P(\mathrm{opt})$ & CZ & ESP & $P(\mathrm{opt})$ & CZ & ESP\\
\midrule
1 & 0.160 & 205 & 0.52 & 0.235 & 529 & 0.19\\
2 & 0.403 & 282 & 0.41 & 0.541 & 941 & 0.05\\
3 & 0.534 & -- & -- & 0.851 & -- & --\\
\bottomrule
\end{tabular}
\end{table}
The XY mixer keeps every sample at Hamming weight $k$, so it falls under case (b) of \cref{thm:qaoa} and gives feasible samples, but the theorem says nothing about its speed of convergence. At comparable or smaller CZ cost it is better here ($0.40$ at $282$ CZ against $0.24$ at $529$ CZ), and at equal depth the Grover mixer is better (for $p=3$, $0.85$ against $0.53$). XY is a local mixer, so it works for counts and not for exact cover or permutations, where single moves do not connect the feasible set.

\paragraph{Hardware-ready example.} The notebook \path{notebooks/General_constraints_Grover_QAOA_hardware.ipynb} is designed to run a six-variable exact-cover instance ($8$ qubits with the registers; $24$ CZ for $A_q$, and $333$ CZ for one amplification iteration on the \texttt{ibm\_fez} layout with the default multi-controlled synthesis used by the script; ideal P(optimum) $0.25\to1$). With the Toffoli V-chain of the depth paragraph above the same iteration needs $206$--$243$ CZ (median $224$, estimated success probability $0.47$ against $0.31$), and the figures below use the default synthesis unless stated. Under the \texttt{ibm\_fez} noise model, $A_q$ alone yields $89$\,\% feasible samples (uniform: $4.7$\,\%) and P(optimum) $0.21$; one iteration yields P(optimum) $0.32$ and $46$\,\% feasible samples. Compilation on the \texttt{ibm\_fez} topology (FakeFez snapshot, $20$ transpiler seeds, \path{experiments/hardware_depth_check.py}) gives legal circuits in every case: $A_q$ uses $22$--$24$ CZ, depth $78$--$86$ and $4.3\,\mu$s on $8$ qubits; one iteration uses $307$--$343$ CZ, depth $768$--$843$ and $30\,\mu$s on $13$ qubits, about a fifth of the median $T_1$ ($144\,\mu$s) and a third of $T_2$ ($98\,\mu$s), with an estimated success probability (product of gate errors) of $0.31$ against $0.90$ for $A_q$ alone. The 333-CZ circuit is therefore predicted to be at the edge of current devices, and most of its ideal gain is lost to noise. We report no data from real hardware. The script \path{experiments/hardware_general.py} runs the circuits with a uniform-state reference, a depth-matched noise control ($A_q$ followed by barrier-separated pairs of identical CZ gates) and several layouts, and \path{experiments/hardware_general_report.py} reports Wilson intervals and three pre-specified one-sided tests ($A_q$ against the uniform state, one iteration against $A_q$, one iteration against the noise control).
